\section{Tarski's undefinability theorem}
\label{sec:second-incompleteness-04}

\begin{definition}[Arithmetical truth definition]
An \(\Lang_A\)-formula \(\mathsf{Tr}(x)\) is an \emph{arithmetical truth
definition} if, for every \(\Lang_A\)-sentence \(\sigma\),
\[
  \mathcal N\entails
  \mathsf{Tr}(\overline{\ulcorner\sigma\urcorner})
  \leftrightarrow\sigma.
\]
The biconditional is evaluated in the standard structure and ranges
metalinguistically over all arithmetic sentences.
\end{definition}

\begin{theorem}[Tarski undefinability]
\label{thm:tarski-undefinability}
There is no arithmetical truth definition for the full set of true
\(\Lang_A\)-sentences.
\textit{Status: complete proof.}
\end{theorem}

\begin{proof}
Suppose \(\mathsf{Tr}(x)\) were such a formula.  Apply the diagonal lemma to
\(\neg\mathsf{Tr}(x)\).  There is a sentence \(L\) such that
\[
  \mathsf Q\derives
  L\leftrightarrow
  \neg\mathsf{Tr}(\overline{\ulcorner L\urcorner}).
\]
Because \(\mathcal N\entails\mathsf Q\), the biconditional is true in
\(\mathcal N\).  But the assumed truth definition also gives
\[
  \mathcal N\entails
  \mathsf{Tr}(\overline{\ulcorner L\urcorner})\leftrightarrow L.
\]
Combining the two yields
\(\mathcal N\entails L\leftrightarrow\neg L\), impossible in the classical
standard structure.
\end{proof}

\begin{corollary}
The set of G\"odel numbers
\[
  \set{\ulcorner\sigma\urcorner:
  \sigma\text{ is an arithmetic sentence and }
  \mathcal N\entails\sigma}
\]
is not definable in \(\mathcal N\) by any \(\Lang_A\)-formula.  In
particular, it is not computably
enumerable.
\end{corollary}

\begin{proof}
A formula defining the displayed set in \(\mathcal N\) would be an
arithmetical truth
definition.  Every computably enumerable set is \(\Sigma^0_1\)-definable, so
nondefinability implies that the set is not computably enumerable.
\end{proof}

\begin{remark}[Partial truth definitions]
For each fixed level of the arithmetical hierarchy, arithmetic can define a
satisfaction predicate for formulas of that bounded complexity.  Tarski's
theorem excludes one arithmetic formula correct for all arithmetic sentences
at once.  This distinction is essential in reflection principles and proof
theory \cite{Tarski1983}.
\end{remark}

